\documentclass[11pt,a4paper]{article}

% --- Pacchetti Matematici Fondamentali ---
\usepackage{amsmath}   % Strutture matematiche avanzate (es. align)
\usepackage{amssymb}   % Simboli matematici e font speciali (es. \mathbb pour i reali R)
\usepackage{amsfonts}  % Font matematici aggiuntivi
\usepackage{amsthm}    % Gestione di teoremi, definizioni e dimostrazioni

% --- Configurazione dei Teoremi ---
\newtheorem{theorem}{Teorema}
\newtheorem{lemma}[theorem]{Lemma}
\newtheorem{proposition}[theorem]{Proposizione}
\theoremstyle{definition}
\newtheorem{definition}[theorem]{Definizione}
\newtheorem{example}[theorem]{Esempio}
\theoremstyle{remark}
\newtheorem{remark}[theorem]{Nota}

\begin{document}

Here we collect some basic facts relating independence of random variables with their characteristic functions. 
The main result is the following, whose proof requires two more theorems proved below.

\begin{theorem}

	Let $X = (X_1,..,X_n)$ and suppose that $\forall \xi \in \mathbb{R}^n \ \phi_X(\xi) = \prod_{j=1,..,n} \phi_{X_j}(\xi_j)$.
	Then $X_1,..,X_n$ are independent.

\end{theorem}

\begin{proof}

	Let $Y = (Y_1,..,Y_n)$ vector of independent random variables, with $Y_j \sim X_j$ for $j = 1,..,n$.
	Now we prove $\phi_X = \phi_Y$, which will give us the thesis applying the unicity theorem.
	\begin{align*}
		\phi_Y(\xi) =& \mathbb{E}[exp(i \cdot \sum_{j=1,..,n} \xi_j Y_j)] = \mathbb{E}[\prod_{j=1,..,n} exp(i \xi_j Y_j)] \\
			=& \prod_{j=1,..,n} \mathbb{E}[exp(i \xi_j Y_j)] = \prod_{j=1,..,n} \phi_{Y_j}(\xi_j) = \prod_{j=1,..,n} \phi_{X_j}(\xi_j)
	\end{align*}

	By the unicity theorem we get \footnote{for any $d$-dimensional random varibale $X$, define 
	$F_X: \mathbb{R}^d \to [0,1], (x_1,..,x_d) \mapsto \mathbb{P}(X_1 \leq x_1,..,X_d \leq x_d)$}. 
	$F_X = F_Y$, so in particular $X_j$'s must be independent since the $Y_j$'s are.

\end{proof}


\begin{theorem}

	Let $X,Y: (\Omega, \mathcal{A}) \to (\mathbb{R}^d, \mathcal{B}(\mathbb{R}^d))$. Suppose $\phi_X = \phi_Y$.
	Then $F_X = F_Y$.

\end{theorem}


\begin{proof}

	$\phi_X = \phi_Y$ means that $\forall \xi \in \mathbb{R}^d \ \mathbb{E}[exp(i \xi X)] = \mathbb{E}[i \xi Y]$.
	Let $p(x) := \sum_{k=1,..,m} c_k exp(i \xi_k x)$, where $c_k \in \mathbb{R}$ and $x,\xi_k \in \mathbb{R}^d$. 
	Such an expression is called a trigonometric polinomial.
	We denote $TP$ the set of all of them.
	By linearity of $\mathbb{E}$ we have $\mathbb{E}[p(X)] = \mathbb{E}[p(Y)]$.
	Consider $C_c(\mathbb{R}^d) := \{f: \mathbb{R}^d \to \mathbb{R} : f \text{ continuous with compact support} \}$.
	Observe that $TP$ has the following properties (meaning that it satisfies Stone-Weierstrass theorem):
	\begin{itemize}
		\item	$TP$ algebra:
		\item	$TP$ separates: given $x \not =y \in \mathbb{R}^d$ one can find $f \in TP$ such that $f(x) \not = f(y)$.
		\item	$TP$ doesn't annihilate points: given $x \in \mathbb{R}^d$, consider any constant, non zero function, $c$. 
			Then $c \in TP$ and $c(x) \not = 0$.
		\item	$TP$ closed under conjugate: $\overline{p}(x) = p(-x)$.
	\end{itemize}

	By Stone-Weierstrass theorem we have $TP$ is dense in $C_c(\mathbb{R}^d)$ with the topology induced by the norm $\|\cdot \|_{\infty}$.
	Equivalently:
	\[
		\forall g \in C_c(\mathbb{R}^d) \forall \epsilon > 0 \exists P_{\epsilon} \in TP \ 
		sup_{x \in \mathbb{R}^d} |g(x) - P_{\epsilon}(x)| < \epsilon
	\]

	Then:
	\[
		\| \mathbb{E}[gX] - \mathbb{E}[gY] \| \leq
		\| \mathbb{E}[gX] - \mathbb{E}[P_{\epsilon}X] \| + \| \mathbb{E}[P_{\epsilon}X] - \mathbb{E}[P_{\epsilon}Y] \| 
		+ \| \mathbb{E}[gY] - \mathbb{E}[P_{\epsilon}Y]\| <
		\epsilon + 0 + \epsilon < 2 \cdot \epsilon
	\]

	Then $\| \mathbb{E}[gX] - \mathbb{E}[gY] \| < 2 \cdot \epsilon$.
	Since $\epsilon$ is arbitrary we get $\forall g \in C_c(\mathbb{R}^d) \ \mathbb{E}[gX] = \mathbb{E}[gY]$.
	Apply now Riesz theorem to get $\mu_X,\mu_Y$ probability measures on $\mathbb{R}^d$ of $X,Y$ respectively. Then we have:
	\[
		\int_{\mathbb{R}^d} g(x) \mu_X(dx) = \int_{\mathbb{R}^d} g(x) \mu_Y(dx)
	\]

	and again by Riesz theorem we have $\mu_X = \mu_y$, from which $F_X = F_Y$ follows.

\end{proof}

\begin{theorem}

	Let $K$ be compact topoplogical space. Let $\mathcal{A}$ subalgebra of $C(K, \mathbb{C})$.
	Suppose that $\mathcal{A}$ has the following properties:
	\begin{itemize}
		\item	separates points: $\forall x \not = y \in K \exists f \in \mathcal{A} \ f(x) \not = f(y)$.
		\item	never annihilites points: $\forall x \in K \exists f \in \mathcal{A} f(x) \not = 0$.
		\item	closed under conjugate: $\forall f \in \mathcal{A} \ \overline{f} \in \mathcal{A}$.
	\end{itemize}
	Then $\mathcal{A}$ dense in $C(K, \mathcal{C})$ with respect to uniform norm.

\end{theorem}

\begin{proof}

	Consider $(P_n)_{n<\omega}$ defined by: $P_0 = c_0, P_{n+1}(t) = P_n(t) + \frac{t^2 - P_n(t)^2}{2}$.
	By induction on $n$ one can prove that, for all $t \in [-1,1]$, $P_n(t)$ is monotone increasing, bounded above by $t \mapsto |t|$.
	Then we can apply Monotone Convergence to get the punctual limit for $(P_n)_{n<\omega}$. Call it $L$.
	It follows that $L(t) = L(t) + \frac{t^2 - L(t)^2}{2}$.
	Then $L(t) = t^2$, from which we deduce (since $L \geq 0$, being $P_n \geq 0$ for all $n < \omega$) that $L(t) = |t|$.
	By Dini theorem, convergence is uniform.
	This means that we just found that $t \mapsto |t|$ can be approximated uniformly by polinomials.
	Now let $\overline{\mathcal{A}}$ be the closure of $\mathcal{A}$ with respect to $\| \cdot \|_{\infty}$.
	We now want to prove that $\overline{\mathcal{A}}$ is a lattice, or equivalently that 
	$max(f,g), min(f,g) \in \overline{\mathcal{A}}$ for all $f,g \in \overline{\mathcal{A}}$.
	Observe that:
	$max(f,g) = \frac{f + g + |f - g|}{2}$ and $min(f,g) = \frac{f + g - |f - g|}{2}$.
	Since $\overline{\mathcal{A}}$ is a vector space, we have $f + g, f - g \in \overline{\mathcal{A}}$.
	We need to prove that if $h \in \overline{\mathcal{A}}$ then $|h| \in \overline{\mathcal{A}}$.
	Observe that, if $h \in C_c(\mathbb{R}^d)$ then $h$ bounded. Let $M$ be such that $\forall x \in K \ |h(x)| \leq M$.
	Let $p(t) = \sum_{k} c_k t^k$ be the polinomial that approximates uniformly $|t|$ in $[-M,M]$. 
	Since $\overline{\mathcal{A}}$ algebra and $h \in \overline{\mathcal{A}}$ we get $\forall k<\omega h^k \in \overline{\mathcal{A}}$.
	Consequently $ph \in \overline{\mathcal{A}}$.
	Since $ph$ uniformly approximates $|h|$, we have $|h| \in \overline{\mathcal{A}}$.
	In conclusion, $\overline{\mathcal{A}}$ is a lattice.
	Now let $f \in C(K), \epsilon > 0$.
	We want to find $g \in \overline{\mathcal{A}}$ such that $\| f - g \| < \epsilon$.
	First consider, for all $x \not = y \in K$, $h_{x,y} \in \mathcal{A}$ such that $h_{x,y}(x) = f(x), h_{x,y}(y) = f(y)$.
	Now fix $y$ and for each $x$ consider $h_{x,y}$.
	Since $h_{x,y} (x) = f(x)$ and $h_{x,y}, f$ continuous, we have that there exists $U_x$ open neighbourhood of $x$ such that 
	$\forall t \in U_x \ h_{x,y}(t) > f(t) - \epsilon$.
	Since $\{U_x\}_{x \in K}$ is an open covering of $K$ and $K$ is compact, we can extract a finite subcovering.
	Let $K \subseteq U_{x_1} \cup .. \cup U_{x_m}$.
	Consider $g_y(t) := max(h_{x_1,y}(t),..,h_{x_m,y}(t))$.
	Since $\overline{\mathcal{A}}$ lattice, we have $g_y \in \overline{\mathcal{A}}$.
	Furthermore the following hold:
	\begin{itemize}
		\item	$\forall t \in K \ g_y(t) > f(t) - \epsilon$
		\item	$g_y(y) = f(y)$
	\end{itemize}

	Since $g_y(y) = f(y)$ we can find (why?? $g_y$ not continuous) $V_y$ open neighbourhood of $y$ such that 
	$\forall t \in V_y \ g_y(t) < f(t) + \epsilon$.
	Again, as $\{V_y\}_{y \in K}$ is an open covering of $K$, compact, we can extract a finite subcovering $V_{y_1},..,V_{y_k}$.
	Define $G(t) := min(g_{y_1}(t),..,g_{y_k}(t))$.
	Since $\overline{\mathcal{A}}$ lattice, $G \in \overline{\mathcal{A}}$.
	Now we have that:
	\begin{itemize}
		\item	$\forall i \ g_{y_i}(t) > f(t) - \epsilon$, so in particular $G(t) > f(t) - \epsilon$.
		\item	$\forall t \in K \ \exists y_i \ t \in V_{y_i}$.
			Since $g_{y_i}(t) < f(t) + \epsilon$, we have $G(t) < f(t) + \epsilon$.
	\end{itemize}
	This means that $\forall t \in K \ f(t) - \epsilon < G(t) < f(t) + \epsilon$.
	Thus, $\| f-G \|_{\infty} < \epsilon$.
	Being $\epsilon$ arbitrary and $G \in \overline{\mathcal{A}}$, we have $f in \overline{\mathcal{A}}$.
	Since $f$ is an arbitrary element of $C(K)$, we conclude that $\overline{\mathcal{A}} = C(K)$.

\end{proof}
\end{document}

